\documentclass[a4paper,11pt,reqno]{amsart}

\usepackage[utf8]{inputenc}
\usepackage[foot]{amsaddr}
\usepackage{amsmath,amsfonts,amssymb,amsthm,mathrsfs,bm}
\usepackage[margin=0.95in]{geometry}
\usepackage{color}
\usepackage[dvipsnames]{xcolor}

\input{toc-config.tex}

\usepackage{mathtools,enumerate,mathrsfs,graphicx}
\usepackage{epstopdf}
\usepackage{hyperref}

\usepackage{latexsym}


\definecolor{CommentGreen}{rgb}{0.0,0.4,0.0}
\definecolor{Background}{rgb}{0.9,1.0,0.85}
\definecolor{lrow}{rgb}{0.914,0.918,0.922}
\definecolor{drow}{rgb}{0.725,0.745,0.769}

\usepackage{listings}
\usepackage{textcomp}
\lstloadlanguages{Matlab}%
\lstset{
    language=Matlab,
    upquote=true, frame=single,
    basicstyle=\small\ttfamily,
    backgroundcolor=\color{Background},
    keywordstyle=[1]\color{blue}\bfseries,
    keywordstyle=[2]\color{purple},
    keywordstyle=[3]\color{black}\bfseries,
    identifierstyle=,
    commentstyle=\usefont{T1}{pcr}{m}{sl}\color{CommentGreen}\small,
    stringstyle=\color{purple},
    showstringspaces=false, tabsize=5,
    morekeywords={properties,methods,classdef},
    morekeywords=[2]{handle},
    morecomment=[l][\color{blue}]{...},
    numbers=none, firstnumber=1,
    numberstyle=\tiny\color{blue},
    stepnumber=1, xleftmargin=10pt, xrightmargin=10pt
}

\numberwithin{equation}{section}
\synctex=1

\hypersetup{
    unicode=false, pdftoolbar=true, 
    pdfmenubar=true, pdffitwindow=false, pdfstartview={FitH}, 
    pdftitle={ELE2024 Coursework}, pdfauthor={A. Author},
    pdfsubject={ELE2024 coursework}, pdfcreator={A. Author},
    pdfproducer={ELE2024}, pdfnewwindow=true,
    colorlinks=true, linkcolor=red,
    citecolor=blue, filecolor=magenta, urlcolor=cyan
}


% CUSTOM COMMANDS
\renewcommand{\Re}{\mathbf{re}}
\renewcommand{\Im}{\mathbf{im}}
\newcommand{\R}{\mathbb{R}}
\newcommand{\N}{\mathbb{N}}
\newcommand{\C}{\mathbb{C}}
\newcommand{\lap}{\mathscr{L}}
\newcommand{\dd}{\mathrm{d}}
\newcommand{\smallmat}[1]{\left[ \begin{smallmatrix}#1 \end{smallmatrix} \right]}

%opening
\title[MATH 2860 (Elementary Differential Equations)]{Homework 5 for MATH 2860 (Elementary Differential Equations)}

\author[Emmanuel Atindama]{E. A. Atindama, PhD Mathematics}

\address[E. A. Atindama]{. Email addresses: \href{emmanuel.atindama@utoledo.edu}{emmanuel.atindama@utoledo.edu} 
% and 
% \href{mailto:a.student@qub.ac.uk}{a.student@qub.ac.uk}.
}
\thanks{Version 0.0.1. Last updated:~\today.}

\begin{document}
\maketitle

Due at 10:00 EST in class

\subsection*{Question Q1}
  A tank initially contains 2,000 litres of water with a concentration of mercury of $6.0\times10^{-5}$ grams
  per litre. Water that contains $3.0\times 10^{-5}$ grams per litre is pumped into the tank at 100 litres per
  hour. The well mixed solution is pumped out of the tank at 100 litres per hour. How long will it
  take for the concentration in the tank to reach $4.5\times 10^{-5}$ grams per litre?


  \begin{center}\setlength{\fboxsep}{10pt}\fcolorbox{yellow!20}{yellow!20}{\parbox{0.9\linewidth}{
    \textbf{Solution}

    Let $A(t)$ be the amount of mercury (in grams) in the tank at time $t$ (hours).  Initial amount
    \[
      A(0)=2000\cdot 6.0\times10^{-5}=0.12\ \text{g}.
    \]
    Inflow brings mercury at rate $100\ \text{L/h}\times 3.0\times10^{-5}\ \text{g/L}=0.003\ \text{g/h}$. The outflow
    rate of mercury is $100\ \text{L/h}\times \dfrac{A(t)}{2000\ \text{L}}=\dfrac{1}{20}A(t)$ (g/h). Hence the ODE is
    \[
      \frac{dA}{dt}=0.003-\frac{1}{20}A.
    \]
    This is linear with steady state $A_{ss}$ satisfying $0.003-\frac{1}{20}A_{ss}=0$, so $A_{ss}=0.003\cdot 20=0.06\ \text{g}$.
    The solution (using integrating factor or standard form) is
    \[
      A(t)=A_{ss} + \big(A(0)-A_{ss}\big)e^{-t/20}
      =0.06 + (0.12-0.06)e^{-t/20}
      =0.06 + 0.06e^{-t/20}.
    \]
    We want the concentration to be $4.5\times10^{-5}\ \text{g/L}$, i.e. amount
    \[
      A_{\text{target}}=2000\cdot 4.5\times10^{-5}=0.09\ \text{g}.
    \]
    Solve $0.09=0.06+0.06e^{-t/20}$:
    \[
      0.03 = 0.06 e^{-t/20} \quad\Rightarrow\quad e^{-t/20}=\tfrac{1}{2}.
    \]
    Hence
    \[
      -\frac{t}{20}=\ln\frac{1}{2} \quad\Rightarrow\quad t=20\ln 2 \approx 13.86\ \text{hours}.
    \]
    Therefore it will take $t=20\ln 2\approx 13.9$ hours for the concentration to reach $4.5\times10^{-5}$ g/L.
    }}
  \end{center}






\subsection*{Question Q2}
  A vat with capacity of 500 gallons initially has 10 lbs. of salt dissolved in 100 gallons of brine. A
  brine solution with 0.5 lbs. of salt per gallon is pumped into the tank at a rate of 4 gallons per
  minute, and the well mixed solution is pumped out at a rate of 2 gallons per minute. How much
  salt will be in the vat when the vat is full? (Hint: first find the solution for the volume of the fluid in the tank, before tackling the amount of salt in it.)

  \begin{center}\setlength{\fboxsep}{10pt}\fcolorbox{yellow!20}{yellow!20}{\parbox{0.9\linewidth}{
    \textbf{Solution}

    Let $V(t)$ be the volume (gallons) in the tank at time $t$ (minutes). With inflow $4$ gal/min and outflow $2$ gal/min,
    \[
      V(t)=100+(4-2)t=100+2t.
    \]
    The tank fills (reaches 500 gal) at time $t_f$ with $100+2t_f=500$, so $t_f=200$ minutes.

    Let $S(t)$ be the amount of salt (lbs) in the tank at time $t$.  The inflow brings salt at rate
    $4\ \text{gal/min}\times 0.5\ \text{lb/gal}=2\ \text{lb/min}$. The outflow removes salt at rate
    $2\ \text{gal/min}\times \dfrac{S(t)}{V(t)} = \dfrac{2S(t)}{100+2t}$ (lb/min). Thus
    \[
      \frac{dS}{dt}=2-\frac{2S}{100+2t}.
    \]
    Simplify the coefficient: $\dfrac{2}{100+2t}=\dfrac{1}{50+t}$, so the ODE becomes
    \[
      \frac{dS}{dt} + \frac{1}{50+t}S = 2.
    \]
    Integrating factor $\mu(t)=\exp\!\big(\int \frac{1}{50+t}\,dt\big)=50+t$. Multiply through:
    \[
      \frac{d}{dt}\big[(50+t)S\big] = 2(50+t).
    \]
    Integrate:
    \[
      (50+t)S = \int 2(50+t)\,dt = 100t + t^2 + C.
    \]
    Use $S(0)=10$ to get $50\cdot 10 = C$, so $C=500$. Therefore
    \[
      S(t)=\frac{t^2+100t+500}{50+t}.
    \]
    When the vat is full at $t=200$ minutes,
    \[
      S(200)=\frac{200^2+100\cdot200+500}{50+200}=\frac{40000+20000+500}{250}=\frac{60500}{250}=242\ \text{lbs}.
    \]
    Thus the vat contains $\boxed{242\ \text{lbs}}$ of salt when full (after 200 minutes).
    }}
    \end{center}



\subsection*{Question Q3}
  A car's fuel tank holds 15 gallons of gasoline. When the tank is full, the gasoline contains 0.015 gallons of ethanol. 
  Suppose that pure gasoline is pumped into the tank at a rate of 1 gallon per minute, and the well mixed solution is 
  pumped out at a rate of 1 gallon per minute. How long will it take for the amount of ethanol in the tank to be 0.003 gallons?

  \begin{center}\setlength{\fboxsep}{10pt}\fcolorbox{yellow!20}{yellow!20}{\parbox{0.9\linewidth}{
    \textbf{Solution}

    Let $E(t)$ be the amount (gallons) of ethanol in the tank at time $t$ (minutes). Volume is constant at 15 gallons.
    Pure gasoline (no ethanol) flows in, so inflow adds no ethanol. Outflow removes ethanol at rate
    $1\ \text{gal/min}\times \dfrac{E(t)}{15}=\dfrac{1}{15}E(t)$. Thus
    \[
      \frac{dE}{dt} = -\frac{1}{15}E.
    \]
    Solve: $E(t)=E(0)e^{-t/15}$ with $E(0)=0.015$. We want $E(t)=0.003$, so
    \[
      0.003 = 0.015 e^{-t/15} \quad\Rightarrow\quad e^{-t/15} = \frac{0.003}{0.015} = 0.2.
    \]
    Hence
    \[
      -\frac{t}{15} = \ln(0.2) \quad\Rightarrow\quad t = -15\ln(0.2) = 15\ln 5 \approx 24.14\ \text{minutes}.
    \]
    Therefore it takes $t=15\ln 5\approx 24.1$ minutes for the ethanol amount to drop to 0.003 gallons.
    
    }}
    \end{center}

\end{document}
